% Integration fragment: no documentclass and no bibliography environment.
% Requires amsmath, amssymb, amsthm, hyperref, and a lemma environment.
% Source numbering below refers to arXiv:2109.02851v2 and
% arXiv:2309.08522v1. This is a derived interface, not a verbatim theorem.

\subsection{The interface for a fixed prime prefix}
\label{sec:prefix-interface}

There are two different weights in this argument. The weight used to
sieve $\mathcal A_b$ is a weight on the remaining divisor $q$, of level
approximately $D/b$. It is not identified with the restriction to
$d=bq$ of a single global weight $\widetilde\lambda_D^{\pm}(d)$.
In particular, the signs and the prefix conditions in the ordinary
Rosser weights change when a prime is removed. For example, for
$q\mid P(z)$ and $z<p$, the relation
$-\lambda_D^-(pq)=\lambda_{D/p}^+(q)$ requires $p^2<D$;
the opposite-sign relation involves $p^3<D$.
For two prefix primes a corresponding lower-weight restriction
can require $p_1p_2^3<D$, which is stronger than
$p_1p_2^2<D$. Thus the assertion below is derived at the level of
the individual box atoms in the sieve construction.

For precise references write $L_{21}$ for
\href{https://arxiv.org/pdf/2109.02851v2}{Lichtman,
\emph{A modification of the linear sieve, and the count of twin primes}},
and $L_{23}$ for
\href{https://arxiv.org/pdf/2309.08522v1}{Lichtman,
\emph{Primes in arithmetic progressions to large moduli, and Goldbach
beyond the square-root barrier}}.
The assertion below is a consequence of their proofs, with the
additional convolution argument supplied here; it is not the literal
statement of $L_{23}$, Proposition~6.6.

Put $\sigma=7/32$ and
\begin{equation}\label{eq:prefix-v}
 v(t)=\min\left\{\frac{1+t}{2},
       \frac{1-(3/2-2\sigma)t}{1+\sigma},
       1-\frac{\sigma+3t}{2}\right\}.
\end{equation}
For $N\in\mathbb N$ let
\[
 E_N(d)=\pi(N;d,N)-\frac{\pi(N)}{\varphi(d)}
 \qquad ((d,N)=1).
\]

\begin{lemma}[Derived prime-prefix sieve interface]
\label{lem:prefix-interface}
Fix $r\in\{1,2\}$ and positive constants $\delta,h,u,t_1,\ldots,t_r$,
and let
\[
 I_i=(N^{t_i},N^{t_i+h}],\qquad
 T_+=\sum_{i=1}^r(t_i+h),\qquad z=N^u.
\]
Assume that the intervals are ordered and disjoint, with every member
of $I_i$ larger than every member of $I_{i+1}$, and $u<t_r$.
Let $\nu$ satisfy
\begin{align}
 \nu+20\delta&\le\inf_{t\in[t_1,t_1+h]}v(t),
       \label{eq:prefix-level-margin}\\
 \sum_{i<j}(t_i+h)+2(t_j+h)&\le\nu-20\delta
       \quad(1\le j\le r),
       \label{eq:prefix-vector-margin}\\
 \nu-T_+-20\delta&>0.
       \label{eq:prefix-residual-positive}
\end{align}
Set $R=N^{\nu-T_+-20\delta}$. Choose the standard modified
linear-sieve weights $\rho_R^{\pm}$ on the remaining divisor $q$,
with sieve precision sufficiently small in terms of the fixed margins.
Use weights independent of the particular primes inside each box.
Then, for every $A>0$,
\begin{equation}\label{eq:prefix-remainder}
 \left|\sum_{\substack{p_i\in I_i\ (1\le i\le r)\\
                         b=p_1\cdots p_r, (b,N)=1}}
       \sum_{\substack{q\mid P(z)\\(q,N)=1}}
       \rho_R^{\pm}(q)E_N(bq)\right|
 \ll_{A,\delta,h,u,\boldsymbol t}N(\log N)^{-A}.
\end{equation}
The constant is independent of the even integer $N$.
The weights can be chosen with the usual linear-sieve main terms
$F(\log R/\log z)$ and $f(\log R/\log z)$, up to the error tending
to zero with the sieve precision. Here it suffices to use the lemma
when $\log R/\log z\ge2$; a zero lower bound requires no distribution
estimate.
\end{lemma}

\begin{proof}
We keep $\delta,h,u,\boldsymbol t$ fixed while $N$ tends to infinity.
Choose the construction precision $\eta>0$ and the associated exponent
mesh sufficiently small compared with $\delta$ and $u$.
All losses below are fixed powers $N^{O(\eta)}$, not factors that are
silently absorbed into an asymptotic error with $\eta$ fixed.
They are paid for by the explicit margins in
\eqref{eq:prefix-level-margin}--\eqref{eq:prefix-residual-positive}.

We use the construction in $L_{21}$, equations~(5.14)--(5.17),
with residual level $R$, rather than imposing a restriction on a
global weight of level $N^\nu$. Its remainder is a finite linear
combination of atoms of the form
\begin{equation}\label{eq:prefix-atom}
 \chi_{J_1}*\cdots*\chi_{J_m}*\beta,
\end{equation}
where $\chi_J$ denotes the prime indicator of a construction interval,
all the displayed primes exceed a cutoff $w$, the small-prime weight
$\beta$ is supported on $k\mid P(w)$ and $k\le R^{O(\eta)}$, and
$|\beta(k)|\le1$. The cutoff is a fixed positive power of $N$,
depending on $\eta$. The number of atoms and the sum of the absolute
values of their scalar coefficients are bounded in terms of $\eta$.
The construction may shrink the rough-part level before adding the
small-prime weight; this changes the residual exponent by $O(\eta)$.
The vectors of rough-prime intervals belong to
$D^{\mathrm{well}}(R^{1+O(\eta)})$. These assertions follow from
the ordinary modified linear sieve, not from its further modification
using $F^*$ in $L_{21}$.

Convolve each atom with $\chi_{I_1}*\cdots*\chi_{I_r}$.
Because $z<\inf I_r$, the prefix remains the largest $r$ primes.
The first $r$ conditions defining $D^{\mathrm{well}}$ hold by
\eqref{eq:prefix-vector-margin}. For a later rough prime $q_j$,
the corresponding condition follows from
\[
 (p_1\cdots p_r)(q_1\cdots q_{j-1}q_j^2)
 \le N^{T_+}R^{1+O(\eta)}
 \le N^{\nu-19\delta}
\]
after making $\eta$ small enough. Refining or rounding all intervals
to the common exponent mesh costs another $N^{O(\eta)}$.
Consequently every concatenated rough vector is in
$D^{\mathrm{well}}(N^{\nu-c\delta})$ for an absolute $c>0$.
Its largest prime has exponent in $[t_1,t_1+h+O(\eta)]$.

Apply $L_{23}$, Proposition~6.1, to a representative of each such
rough box. The partition of the prime indices furnished by that
proposition is kept fixed throughout that box. The exponent mesh
then enlarges each of the three factor supports by only
$N^{O(\eta)}$; it does not require choosing a partition separately
for every integer in the box. This is the step implemented in
$L_{21}$, Lemma~5.3. Its proof uses the factorization of a rough
box and the size of the small-prime weight, not membership of the
concatenated vector in a global $D^+$ or $D^-$.

For completeness, the treatment of signs, repeated intervals, and
coefficients can be made explicit. Write an ordered-tuple rough atom
as $T_J=\chi_{J_1}*\cdots*\chi_{J_m}$, with the squarefree restriction
on the rough product understood. Retain exactly its scalar
coefficient from the sieve construction, including the original
factorials for repeated interval labels in $L_{21}$, (5.10).
For an explicit convention, let $m_I$ be the multiplicity of each
interval label and put $M_J=\prod_I m_I!$. If $A_J$ is the indicator
of the corresponding sorted squarefree prime-box product, then
\begin{equation}\label{eq:box-symmetry-normalization}
 \mathbf1_{\rm squarefree}T_J=M_J A_J.
\end{equation}
Indeed each squarefree product has exactly $m_I!$ assignments to
the slots carrying label $I$. With the tuple-convolution convention
of $L_{21}$, (5.10), the scalar already includes
$(-1)^m/M_J$, in addition to any other original scalar factors.
Keep this reciprocal factorial exactly once; it is not inserted a
second time during the three-group normalization below.
If $\mathcal I_1,\mathcal I_2,\mathcal I_3$ is the fixed partition
of its prime slots (including the prefix), put
\[
 K_i=|\mathcal I_i|!,\qquad
 \gamma_i=K_i^{-1}\mathop{*}_{j\in\mathcal I_i}\chi_{J_j}.
\]
The empty convolution is the point mass at $1$. Unique prime
factorization bounds the number of ordered representations by
$K_i$, even on products containing repeated primes, so
$|\gamma_i(n)|\le1$ for every integer $n$. The unrestricted
convolution identity is exactly
\[
 T_J=(K_1K_2K_3)\,\gamma_1*\gamma_2*\gamma_3.
\]
Absorb the displayed, fixed factorial product into the original
scalar coefficient. The extra nonsquarefree rough products will
be removed by the error estimate below; no claim of a unique
ordered prime-slot representation is needed.
The small-prime factor is combined with the third rough factor.
The factorization of an integer into its primes
at most $w$ and its primes greater than $w$ is unique. Hence, if
$\gamma_3$ is the normalized third rough factor,
\[
 |(\gamma_3*\beta)(n)|
 =|\gamma_3(n_{>w})\beta(n_{\le w})|\le1.
\]
This argument allows either sign of $\beta$ and does not assume
positivity. Its support is enlarged by at most $R^{O(\eta)}$;
the degree of every factor-size inequality in Proposition~6.1 is
bounded, so that the resulting loss is again $N^{O(\eta)}$.

One can also implement the interval normalization using unrestricted
prime convolutions and remove collisions afterwards. This gives
exact convolution identities away from integers divisible by
$p^2$ for a rough prime $p>w$. The omitted discrepancy is at most
\[
 \sum_{\substack{d\le N^{\nu+O(\eta)}\\
                  p^2\mid d\ \mathrm{for\ some}\ p>w}}
   \tau_C(d)\left(\frac{N}{d}+1+\frac{N}{\varphi(d)}\right)
 \ll_C (N+N^{\nu+O(\eta)})w^{-1}(\log N)^{C'},
\]
where $C,C'$ are fixed with the precision. Since $w$ is a fixed
positive power of $N$ and $\nu<1$, this is
$O_A(N(\log N)^{-A})$ for every fixed $A$. Thus the squarefree
convention in the sieve causes no unaccounted convolution error.

We spell out the source of uniformity in the residue, since the
displayed constant in $L_{23}$, Proposition~6.6, is written with
a subscript $a$. The input used for a box atom is $L_{23}$,
Proposition~5.2, not an inference from that subscript. Proposition~5.2
has the explicit supremum $0<|a|<x^{1+\epsilon}$ and permits arbitrary
complex, $1$-bounded coefficients on its three modulus factors.
With its Type~II length denoted here by $U$, its size conditions are
\begin{align*}
 Q_1&<Ux^{-\epsilon},&
 U^2Q_2Q_3^2&<x^{1-8\epsilon},\\
 |a|^\sigma UQ_1Q_2^5Q_3^2&<x^{2-15\epsilon},&
 (UQ_3/Q_2)^{2\sigma}UQ_1Q_2^5Q_3^2
     &<x^{2-15\epsilon}.
\end{align*}
Here is a quantitative check of the parameters, including the
endpoints of the Type~II range. Fix a factorization margin
$0<\zeta<10^{-5}$, sufficiently small in terms of $\delta$, and take
$\eta\ll\epsilon\ll\zeta$. For
\[
 x^{2\zeta}\le U\le x^{1/3+\zeta/2},\qquad A_0=Ux^{-\zeta},
\]
we have $x^\zeta\le A_0\le x^{1/3-\zeta/2}$, exactly the allowed
range of the parameter $A$ in Proposition~6.1. Denote its three
factors by $a_0,b_0,c_0$. Apart from $x^{O(\eta)}$ box enlargements,
put $Q_1=a_0$, $Q_2=b_0$, and $Q_3=c_0$ and then split into dyadic
intervals. In particular $Q_1\le A_0<Ux^{-\epsilon}$.
Replacing $A_0$ by $U$ in the three monomials in (6.3) costs,
respectively,
\[
 x^{2\zeta},\qquad x^\zeta,\qquad
 x^{(1+2\sigma)\zeta}.
\]
Their original savings are $3\zeta,3\zeta,2\zeta$.
Thus the smallest remaining saving is
$(1-2\sigma)\zeta=9\zeta/16>0$.
Choosing, for example, $\epsilon<\zeta/1000$ and then the construction
mesh sufficiently small makes all four displayed conditions strict.
If a dyadic value of $Q_2$ is smaller than its support endpoint,
the last monomial is still bounded by its endpoint value: its power
of $Q_2$ is $5-2\sigma>0$. This justifies the same conclusion on
every dyadic subbox.

There are two additional bookkeeping points in passing to primes.
First, the primes in $\pi(N)$ cannot all be treated on the artificial
scale $x=N$. Remove the initial range $p\le N^{1-\kappa}$, with a
fixed $0<\kappa\ll\epsilon$. Its progression term is bounded by
\[
 \sum_{p\le N^{1-\kappa}}\tau_C(N-p)
       \ll N^{1-\kappa+o(1)},
\]
and its mean term is $O(N^{1-\kappa}(\log N)^{C'})$.
On the remaining dyadic intervals $[x,2x]$, use
$\alpha=\log N/\log x=1+O(\kappa)$ in (6.3).
The exponent coordinates, total level, and the function
$v_\alpha(t)$ from (6.2) differ from their $N$-scale counterparts
by $O(\kappa)$. These changes are paid for by the fixed margins.
Also $N<x^{1+\epsilon}$ after making $\kappa$ small enough.
Thus Proposition~5.2 applies with residue $a=N$ on every interval;
it is not enough merely to set $\alpha=1$ on the first interval.
If the full support is at most $N^{1/2-c\delta}$, ordinary
Bombieri--Vinogradov suffices for each bounded atom. Only the
remaining case needs Proposition~6.1, whose section assumes its
level parameter is at least $1/2$.

Second, apply the $k=3$ Heath--Brown identity of $L_{23}$,
Lemma~3.3, as in the proof of Theorem~1.7, but with middle interval
$[x^{2\zeta},x^{1/3+\zeta/2}]$. The six dyadic variables are
$N_1,N_2,N_3,M_1,M_2,M_3$, with $M_i\le(2x)^{1/3}$.
If a variable lies in that interval, take it as the Type~II variable
$U$ and combine the other five. This uses exactly the coefficient
sequences of the cited Heath--Brown argument, rather than a new
Siegel--Walfisz assertion for an arbitrary convolution. The chosen
sequence is $1$ or $\mu$, with the roughness cutoff
$P^-(n)>x^{1/(\log\log x)^3}$ used there; the logarithm is handled
by partial summation. Lemma~\ref{lem:rough-sw} verifies the full
Siegel--Walfisz condition~(2.2), including its auxiliary coprimality
and the rough cutoff. Divisor-bounded
coefficients on the other factor are reduced to $1$-bounded ones
by the truncation argument in the proof of Proposition~3.1, using
Lemmas~4.1 and~4.3; the resulting powers of $\log x$ are absorbed
by increasing the arbitrary exponent in the error estimate.

If no variable lies in the middle interval, every $M_i$ is less
than $x^{2\zeta}$, and at most two of the $N_i$ exceed
$x^{1/3+\zeta/2}$. If just one does, the product of the other five
is at most $x^{10\zeta}$, so the long variable is at least
$x^{1-10\zeta+o(1)}$ and Lemma~3.4 applies. If two do, the
product of the remaining four is at most $x^{8\zeta}$; apply
Proposition~3.2, which has an explicit supremum over all $a\ne0$.
The total modulus support is at most $x^{3/5}$, after making the
fixed margins sufficiently small: indeed $\max_{t\ge0}v(t)
=153/256<3/5$. For small $\zeta,\epsilon$ this is less than
$(x/x^{8\zeta})^{2/3-3\epsilon}$, as required by that proposition,
and less than the long variable times $x^{-\epsilon}$ as required
by Lemma~3.4. Having no large variable is impossible because the
product of all six variables is comparable with $x$.
These arguments also allow the sharp product cutoff on a dyadic
interval, by the same separation and partial summation used in
the proof of Proposition~3.1. Restrictions $(d,N)=1$ split into
$(q_i,N)=1$ on the modulus factors and preserve their bounds.
No fixed-residue estimate from $L_{21}$, Theorem~2.5, is used here.

We obtain a remainder $O_A(N(\log N)^{-A})$ for each atom, uniformly
in $N$. Summing the fixed number of normalized atoms proves
\eqref{eq:prefix-remainder}. Notice that the signs of the atoms were
never used to deduce a distribution estimate. They only enter when
reassembling the residual upper or lower sieve weights.

Finally, the main term is evaluated before convolving in the prefix.
For each $b=p_1\cdots p_r$ the residual density is
$g_N(\ell)=1/(\ell-1)$ for $\ell\nmid N$ and $g_N(\ell)=0$
for $\ell\mid N$, because $(b,q)=1$. Uniformity in $N$ here
requires a minor precaution: the two-sided condition (2.3) of
$L_{21}$ need not hold with an absolute constant all the way down
to $2$, since $N$ can contain many small primes.
Use the small-prime presieve already present in (5.15), with
$w=R^{\eta^2}$ and small-factor level $R^{O(\eta)}$.
Below $w$ the one-sided dimension-one estimate is uniform, since
deleting prime factors from the usual product only decreases its
ratio. The fundamental lemma therefore gives the full factor
$V_N(w)$, with a relative error tending to zero with $\eta$.
Above $w$ the two-sided estimate is uniform: for $w\le y<z$,
\[
 \sum_{\substack{\ell\mid N\\\ell\ge w}}
       \frac1{\ell-1}
 \ll\frac{\log N}{w\log w},\qquad
 \frac{V_N(y)}{V_N(z)}
 =\frac{\log z}{\log y}
   \left(1+O\!\left(\frac1{\log y}
                    +\frac{\log N}{w\log w}\right)\right).
\]
Consequently the usual rough-box evaluation in Theorem~2.10
applies with constants independent of $N$. Reattaching the
presieve gives the usual $F(s),f(s)$, with
$s=\log R/\log z$, and an error tending to zero with the precision.
The entire product $V_N(z)$ is retained; it is not replaced by a
product with the prime divisors of $N$ restored. On summing over
the prefix box the main factor is
\[
 \pi(N)V_N(z)
 \sum_{\substack{p_i\in I_i\ (1\le i\le r)\\(p_1\cdots p_r,N)=1}}
       \frac1{\varphi(p_1\cdots p_r)},
 \qquad
 V_N(z)=\prod_{\substack{\ell<z\\\ell\nmid N}}
       \left(1-\frac1{\ell-1}\right).
\]
Its upper and lower multipliers are respectively
$F(s)+O(\eta)$ and $f(s)+O(\eta)$, with the stated precision
error interpreted on a fixed compact range of $s\ge2$.
This evaluation uses weights constructed for the residual sequence;
it does not replace their main term by a global-prefix weight sum.
\end{proof}

\paragraph{Order of limits and use in the integrals.}
For each fixed positive level margin $\delta$, first choose a finite
exponent partition and sieve precision so that all box enlargements
are smaller than that margin, and then let $N\to\infty$.
Next let the partition width and sieve precision tend to zero,
still within the fixed margin, and finally let $\delta\downarrow0$.
Continuity of $F,f$ on the relevant compact intervals gives the
arguments $(v(t_1)-t_1)/u$ and
$(v(t_1)-t_1-t_2)/u$. Prefix boxes meeting ordering boundaries
can be discarded for lower bounds, and bounded with a classical
sieve for upper bounds; their main contribution tends to zero with
the partition width. A boundary on which the prefix
$D^{\mathrm{well}}$ condition is an equality is treated by a thin
strip before taking the limit, not by applying the lemma at equality.

\paragraph{The middle branch in the level.}
The two-branch form of $v$ used in the numerical calculation is exact.
Put $A(t)=(1+t)/2$, $B(t)=(32-34t)/39$, and
$C(t)=57/64-3t/2$. Then
\[
 B-A=\frac{25-107t}{78},\qquad
 C-A=\frac{25}{64}-2t,\qquad
 B-C=\frac{1568t-175}{2496}.
\]
For $t\le25/128$ the minimum is $A$, whereas for $t\ge25/128$
the minimum is $C$ (the zero of $B-C$ is $25/224<25/128$).
Thus for every $t\ge0$,
\[
 v(t)=\begin{cases}
 (1+t)/2,&t\le25/128,\\
 57/64-3t/2,&t\ge25/128.
 \end{cases}
\]

\paragraph{Audit status of this interface.}
This fragment supplies an adaptation of the box-atom proof, rather
than quoting a published prime-prefix theorem. Independent review
should check its use of the modified-sieve atom decomposition and
the transfer of the positive factorization margins through the
Type~II argument. No numerical Chen constant is asserted here, and
the distribution and linear-sieve theorems cited above are external
inputs rather than independently reproved results.
